% Part: lambda-calculus
% Chapter: introduction
% Section: reduction

\documentclass[../../../include/open-logic-section]{subfiles}

\begin{document}

\olfileid[id]{lam}{int}{red}
\olsection{Reduksi Suku Lambda}

Apa yang dapat dilakukan dengan suku lambda? Menyederhanakannya. Jika $M$ dan
$N$ adalah sebarang suku lambda dan $x$ sebarang variabel, kita dapat memakai
$\Subst{M}{N}{x}$ untuk menyatakan hasil penyubstitusian $N$ sebagai pengganti
$x$ dalam~$M$, setelah mengganti nama setiap variabel terikat dalam~$M$ yang
akan mengganggu variabel bebas dalam~$N$ sesudah substitusi. Sebagai contoh,
\[
\Subst{(\lambd[w][xxw])}{yyz}{x} = \lambd[w][(yyz)(yyz)w].
\]

\begin{digress}
Notasi alternatif untuk substitusi adalah $[N/x]M$, $[x/N]M$, dan juga
$M[x/N]$. Berhati-hatilah!
\end{digress}

Secara intuitif, $(\lambd[x][M])N$ dan $\Subst{M}{N}{x}$ mempunyai makna yang
sama; tindakan mengganti suku pertama dengan suku kedua disebut
\emph{kontraksi-$\beta$}. Suku $(\lambd[x][M])N$ disebut \emph{redex}, dan
$\Subst{M}{N}{x}$ disebut \emph{kontraktumnya}. Secara umum, jika suatu suku
$P$ dapat diubah menjadi~$P'$ melalui kontraksi-$\beta$ pada suatu subsuku,
kita mengatakan bahwa $P$ \emph{tereduksi-$\beta$ menjadi $P'$ dalam satu
langkah}, dan menulis $P \redone P'$. Jika dari $P$ kita dapat memperoleh~$P'$
dengan sejumlah reduksi satu-langkah (mungkin nol), maka $P$
\emph{tereduksi-$\beta$} menjadi~$P'$; secara simbolis, $P \red P'$. Suku yang
tidak dapat direduksi-$\beta$ lebih lanjut disebut \emph{tak tereduksi-$\beta$}
atau \emph{normal-$\beta$}. Kita akan mengatakan ``tereduksi'' sebagai ganti
``tereduksi-$\beta$'', dan seterusnya, apabila konteksnya jelas.

Mari kita perhatikan beberapa contoh.
\begin{enumerate}
\item Kita mempunyai
\begin{align*}
(\lambd[x][xxy]) \lambd[z][z] & \redone (\lambd[z][z])(\lambd[z][z]) y \\
& \redone (\lambd[z][z]) y \\
& \redone y.
\end{align*}
\item ``Menyederhanakan'' suatu suku dapat membuatnya lebih rumit:
\begin{align*}
(\lambd[x][xxy])(\lambd[x][xxy]) & \redone (\lambd[x][xxy])(\lambd[x][xxy])y \\
& \redone (\lambd[x][xxy])(\lambd[x][xxy])yy \\
& \redone \dots
\end{align*}
\item Penyederhanaan juga dapat membiarkan suatu suku tidak berubah:
\[
(\lambd[x][xx])(\lambd[x][xx]) \redone (\lambd[x][xx])(\lambd[x][xx]).
\]
\item Selain itu, beberapa suku dapat direduksi melalui lebih dari satu cara;
  misalnya,
\[
(\lambd[x][(\lambd[y][yx]) z)] v \redone (\lambd[y][yv]) z
\]
dengan mengontraksi aplikasi terluar; dan
\[
(\lambd[x][(\lambd[y][yx]) z)] v \redone (\lambd[x][zx]) v
\]
dengan mengontraksi aplikasi terdalam. Namun, perhatikan bahwa dalam kasus ini
kedua suku selanjutnya tereduksi menjadi suku yang sama, yaitu~$zv$.
\end{enumerate}

Hasil akhir pada contoh terakhir bukanlah suatu kebetulan, melainkan
memperlihatkan sifat kalkulus lambda yang mendalam dan penting, yang dikenal
sebagai ``sifat Church--Rosser.''

\end{document}
